\documentclass[a4paper]{article}

\usepackage[polish]{babel}
\usepackage[T1]{fontenc}
\usepackage[utf8]{inputenc}
\usepackage{graphicx}
\usepackage{fancyhdr}
\usepackage{hyperref}
\pagestyle{fancy}

\lhead[]{}
\chead[Fundacja CoderDojo Polska]{Fundacja CoderDojo Polska}
\rhead[]{}
\cfoot[Treść udostępniana na licencji Creative Commons – Uznanie Autorstwa – Na Tych Samych Warunkach]{Treść udostępniana na licencji Creative Commons – Uznanie Autorstwa – Na Tych Samych Warunkach}

\title{Jak wesprzeć działanie Fundacji CoderDojo Polska?}
\author{Fundacja CoderDojo Polska}
\date{\today}


\begin{document}
\maketitle
\includegraphics[width=\textwidth]{logo_CD.pdf}


\section{Dołącz do nas!}
\begin{itemize}
\item Załóż nowe dojo. Każdy nowy klub jest dla nas ważny i daje kolejnym dzieciom możliwość rozwijania swoich pasji. 
\item Zawsze przyda nam się nowy mentor. Jeśli masz czas i ochotę nim zostać – \href{mailto:dojo@coderdojo.org.pl}{napisz do nas}. Znajdziemy Ci zajęcie 
\item Zostań naszą gwiazdą. Jeśli nie masz czasu, żeby przychodzić do nas często – przyjdź raz. Pokaż dzieciom coś fajnego, coś w czym jesteś dobry. 
\item Przyjdź na zajęcia przydaj się na coś: wesprzyj dzieci, pomóż rozstawić sprzęt, upiecz ciasteczka, pokrój marchewkę na przekąski dla dzieci. Każda pomoc jest cenna. 
\end{itemize}

\section{Wesprzyj nas!}
\begin{itemize}
\item Udostępnij nam salę na organizację dojo. 
\item Wypożycz lub oddaj sprzęt, który nie jest Ci już potrzebny. Nie musi to być najnowszy MacBook, ważne, żeby był sprawny i miał działające WiFi. 
\item Jeśli masz firmę, która działa w branży IT opowiedz o nas swoim współ/pracownikom. Zawsze potrzebujemy mentorów do pomocy. 
\item Zasponsoruj dzieciom naklejki, koszulki, fajne gadżety. 
\item Pomóż nam zdobyć wsparcie lokalnych firm lub władz. Każda współpraca ma dla nas potencjał. 
\item Opowiedz o nas lokalnym mediom, jeśli znasz jakichś dziennikarzy. Każda promocja pomaga nam zakładać nowe dojo.
\end{itemize}

\section{Like!}
\begin{itemize}
\item Daj nam \href{https://www.facebook.com/CoderDojoPL/}{lajka}! Udostępnij naszą stronę w mediach społecznościowych, napisz o nas swoim znajomym, opowiedz na spotkaniu w szkole. 
\item Podlinkuj u siebie naszą stronę www, wstaw nasz banner. 
\item Jeśli jesteś dziennikarzem lub blogerem – napisz o nas. Przyjedź, zobacz jak wyglądają zajęcia i zostań naszym przyjacielem. 
\item Powieś u siebie nasz plakat, rozłóż ulotki. 
\item Masz jakieś refleksje, pomysły, uważasz, że coś możemy zrobić lepiej? \href{mailto:dojo@coderdojo.org.pl}{Napisz do nas} – każdego głosu słuchamy uważnie.
\end{itemize}

\end{document}